\section{CONC-002: Async Hook Execution Ordering}
\label{sec:conc-002}

\subsection*{Metadata}
\begin{tabular}{ll}
\textbf{Issue ID} & CONC-002 \\
\textbf{Severity} & Medium \\
\textbf{Category} & Concurrency \\
\textbf{Branch} & \branchref{fix/CONC-002-async-hook-execution-ordering} \\
\textbf{Changes} & \branchdiff{fix/CONC-002-async-hook-execution-ordering} \\
\end{tabular}

\subsection*{Executive Summary}
The \texttt{executeHook} function in \texttt{src/hooks.ts} spawns a child process and races a \texttt{setTimeout} against the \texttt{close} event. If the timeout fires at exactly the same time as the hook exits, both \texttt{reject} and \texttt{promiseResolve} can be called, leading to a double-settle on the promise. Node.js promises ignore the second call, but when combined with the \texttt{setTimeout} callback calling \texttt{child.kill()} after the process has already exited, this creates an unhandled rejection risk.

The bug is in \texttt{src/hooks.ts:92-113}, where the timeout-based rejection and close-based resolution are not properly gated. The fix introduces a \texttt{settled} flag that ensures the promise settles exactly once, regardless of event ordering.

\subsection*{Root Cause Analysis}

The race occurs between the \texttt{setTimeout} callback and the \texttt{close} event handler:

\begin{lstlisting}[language=TypeScript]
// hooks.ts:59-117 — Before fix
const timeout = setTimeout(() => {
    child.kill("SIGTERM");
    reject(new Error(`Hook "${hook.script}" timed out after 10s`));
}, 10_000);

child.on("error", (err) => {
    clearTimeout(timeout);
    reject(new Error(`Hook "${hook.script}" failed to start: ${err.message}`));
});

child.on("close", (code) => {
    clearTimeout(timeout);
    if (code !== 0) {
        reject(new Error(/* ... */));
        return;
    }
    try {
        const result = JSON.parse(stdout.trim());
        promiseResolve(result);
    } catch {
        promiseResolve({ action: "allow" });
    }
});
\end{lstlisting}

If the \texttt{close} event and \texttt{setTimeout} fire in the same microtask tick, both \texttt{reject} and \texttt{promiseResolve} are called. While only one settles the promise, the \texttt{child.kill("SIGTERM")} in the timeout callback may execute after the process has already exited, sending a signal to a dead process (harmless but wasteful).

\subsection*{Fix Implementation}

The fix introduces a \texttt{settled} guard and a unified \texttt{settle} function:

\begin{lstlisting}[language=TypeScript]
// hooks.ts — After fix
async function executeHook(
    hook: HookDefinition,
    agentDir: string,
    input: Record<string, any>,
): Promise<HookResult> {
    return new Promise<HookResult>((promiseResolve, reject) => {
        const child = spawn(/* ... */);
        let stdout = "";
        let stderr = "";
        let settled = false;

        function settle(resolveVal: HookResult | null, rejectErr: Error | null) {
            if (settled) return;
            settled = true;
            clearTimeout(timeout);
            if (rejectErr) reject(rejectErr);
            else if (resolveVal) promiseResolve(resolveVal);
        }

        const timeout = setTimeout(() => {
            child.kill("SIGTERM");
            settle(null, new Error(
                `Hook "${hook.script}" timed out after 10s`));
        }, 10_000);

        child.on("error", (err) => {
            settle(null, new Error(
                `Hook "${hook.script}" failed to start: ${err.message}`));
        });

        child.on("close", (code) => {
            if (code !== 0) {
                settle(null, new Error(
                    `Hook "${hook.script}" exited with code ${code}: ${stderr.trim()}`));
                return;
            }
            try {
                const result = JSON.parse(stdout.trim());
                settle(result, null);
            } catch {
                settle({ action: "allow" }, null);
            }
        });
    });
}
\end{lstlisting}

The key change is that all exit paths---timeout, error, close---route through \texttt{settle()}, which checks the \texttt{settled} flag and clears the timeout. This guarantees the promise resolves or rejects exactly once.

\subsection*{Affected Files}
\begin{itemize}
    \item \srcfile{src/hooks.ts} --- added \texttt{settled} guard to prevent double settle
    \item \srcfile{src/\_\_tests\_\_/hooks-conc.test.ts} --- new concurrency tests
\end{itemize}

\subsection*{Verification}
\begin{lstlisting}[language=TypeScript]
// hooks-conc.test.ts
describe("CONC-002: Hook execution ordering", () => {
    it("executeHook settles exactly once even if close and timeout race", async () => {
        const { executeHook } = await import("../hooks.js");
        const fastHook = {
            script: "echo '{\"action\":\"allow\"}'",
            baseDir: "/tmp",
        } as any;

        const result = await executeHook(fastHook, "/tmp", {});
        assert.ok(result);
        assert.strictEqual(result.action, "allow");
    });

    it("executeHook rejects on timeout", async () => {
        const { executeHook } = await import("../hooks.js");
        const hangingHook = {
            script: "sleep 15",
            baseDir: "/tmp",
        } as any;

        await assert.rejects(
            () => executeHook(hangingHook, "/tmp", {}),
            { message: /timed out/ },
            "Should reject with timeout error",
        );
    });
});
\end{lstlisting}
